RLHF-trained language models are widely criticized for producing verbose, bland, and homogeneous outputs---the ``all AI sounds the same'' problem. A natural hypothesis is that preference aggregation across many annotators is the root cause: averaging diverse human preferences into a single reward signal may compress individual variation into a generic mean. We test this hypothesis by training DPO models on individual simulated annotators with diverse verbosity preferences and comparing them to an aggregate-trained baseline. We first analyze 1,393 users in the \prism dataset and find that while 84.1\% of individuals prefer longer responses, per-user preferences vary significantly ($H{=}920.81$, $p{<}10^{-93}$). We then train DPO on three annotators spanning the verbosity spectrum: a length-loving annotator, a brevity-preferring annotator, and a moderate annotator. Aggregate DPO amplifies verbosity by 21\% over the base model (143.6 vs.\ 118.8 words), while the brevity-preferring annotator reduces it by 15\% (100.8 words, $p{=}0.042$). However, an LLM judge rates distinctiveness at only 1.3--1.6 out of 5 for \emph{all} conditions---base, individual, and aggregate alike. We conclude that RLHF's verbosity problem is largely a preference-aggregation artifact that personalization can mitigate, but the deeper blandness pervading model outputs is rooted in pretraining and supervised fine-tuning, not in RLHF itself.
